% Źródła: Eves s. 287--288, 314--315; Wikipedia: Bernhard Riemann.
% https://en.wikipedia.org/wiki/Bernhard_Riemann

\section{Bernhard Riemann} % lang-pl
\section{Bernhard Riemann} % lang-it
\label{sekcja_riemann}

Hipotezę kąta rozwartego wszyscy dotąd odrzucali, bo przeczyła nieskończoności prostej. % lang-pl
W tej sekcji zobaczymy, jak uczeń Gaussa odróżni nieograniczoność od nieskończoności i jak jeden wykład otworzy drugi okres w dziejach geometrii nieeuklidesowej. % lang-pl
L'ipotesi dell'angolo ottuso l'avevano finora scartata tutti, perché contraddiceva l'infinità della retta. % lang-it
In questa sezione vedremo come un allievo di Gauss distinguerà l'illimitatezza dall'infinità e come una sola conferenza aprirà un secondo periodo nella storia della geometria non euclidea. % lang-it

Bernhard Riemann (1826--1866) urodzi się w Breselenz w Królestwie Hanoweru, będzie studiował w Getyndze u Gaussa, który zostanie jego promotorem i zachęci go do matematyki zamiast teologii. % lang-pl
Zachoruje na gruźlicę i umrze we Włoszech w wieku trzydziestu dziewięciu lat. % lang-pl
Bernhard Riemann (1826--1866) nascerà a Breselenz, nel Regno di Hannover, studierà a Gottinga con Gauss, che diventerà il suo relatore e lo incoraggerà alla matematica anziché alla teologia. % lang-it
Si ammalerà di tubercolosi e morirà in Italia a trentanove anni. % lang-it
\index[persons]{Riemann, Bernhard}%

Wykład habilitacyjny \emph{Über die Hypothesen, welche der Geometrie zu Grunde liegen} wygłosi 10 czerwca 1854 roku w Getyndze; wydrukuje go dopiero w 1868 roku, już po śmierci Riemanna, Richard Dedekind. % lang-pl
Riemann omówi w nim pojęcia nieograniczoności i nieskończoności; gdy różnica między nimi stanie się jasna, da się zrealizować równie niesprzeczną geometrię hipotezy kąta rozwartego, jeśli zmodyfikuje się postulaty Euklidesa \cite[s. 287]{eves1_1972}: % lang-pl
La lezione di abilitazione \emph{Über die Hypothesen, welche der Geometrie zu Grunde liegen} la terrà il 10 giugno 1854 a Gottinga; la stamperà solo nel 1868, già dopo la morte di Riemann, Richard Dedekind. % lang-it
Riemann vi discuterà i concetti di illimitatezza e infinità; una volta chiarita la differenza, sarà possibile realizzare una geometria altrettanto coerente per l'ipotesi dell'angolo ottuso, se si modificano i postulati di Euclide \cite[p. 287]{eves1_1972}: % lang-it
\index[persons]{Dedekind, Richard}%
\begin{itemize}
\item dwa różne punkty wyznaczają co najmniej jedną prostą; % lang-pl
\item prosta jest nieograniczona; % lang-pl
\item każde dwie proste na płaszczyźnie przecinają się. % lang-pl
\item due punti distinti determinano almeno una retta; % lang-it
\item la retta è illimitata; % lang-it
\item due rette qualsiasi di un piano si intersecano. % lang-it
\end{itemize}
Eves zauważa, że pierwsze dwa z tych zdań to po prostu postulaty Euklidesa wzięte dosłownie: Euklides miał na myśli, że dwa punkty wyznaczają \emph{dokładnie} jedną prostą i że jest ona \emph{nieskończona}, ale tego nie napisał \cite[s. 288]{eves1_1972}. % lang-pl
Eves osserva che le prime due sono semplicemente i postulati di Euclide presi alla lettera: Euclide intendeva che due punti determinano \emph{esattamente} una retta e che essa è \emph{infinita}, ma non lo scrisse \cite[p. 288]{eves1_1972}. % lang-it

Wykład Riemanna otworzy drugi okres rozwoju geometrii nieeuklidesowej: zamiast metod elementarnej geometrii syntetycznej stosować się będzie metody geometrii różniczkowej. % lang-pl
Riemann stworzy w ten sposób całą klasę nietradycyjnych geometrii, zwanych \emph{riemannowskimi}, a uogólnienie pojęcia przestrzeni, do którego prowadzi jego wykład, zastosuje się później w teorii względności \cite[s. 288]{eves1_1972}. % lang-pl
Eugenio Beltrami powoła się na ten wykład w swoich pracach o modelach (zobacz sekcję \ref{sekcja_beltrami_cayley_klein}), które ukażą się w tym samym 1868 roku, w którym Dedekind wydrukuje wykład. % lang-pl
La conferenza di Riemann aprirà un secondo periodo nello sviluppo della geometria non euclidea: al posto dei metodi della geometria sintetica elementare si useranno quelli della geometria differenziale. % lang-it
Riemann creerà così un'intera classe di geometrie non tradizionali, dette \emph{riemanniane}, e la generalizzazione del concetto di spazio a cui porta la sua conferenza sarà poi applicata alla teoria della relatività \cite[p. 288]{eves1_1972}. % lang-it
Eugenio Beltrami si richiamerà a questa conferenza nei suoi lavori sui modelli (si veda la sezione \ref{sekcja_beltrami_cayley_klein}), che usciranno nello stesso 1868 in cui Dedekind stamperà la conferenza. % lang-it
\index{geometria!riemannowska} % lang-pl
\index{geometria!riemanniana} % lang-it

Eves wspomina jeszcze inną nietradycyjną geometrię: Max Dehn zbuduje taką, w której zawieszony jest aksjomat Archimedesa (geometria nieArchimedesowa) \cite[s. 288]{eves1_1972}. % lang-pl
Eves cita ancora un'altra geometria non tradizionale: Max Dehn ne costruirà una in cui è sospeso il postulato di Archimede (geometria non archimedea) \cite[p. 288]{eves1_1972}. % lang-it
\index[persons]{Dehn, Max}%

\todofoot{Geometria riemannowska (metryka, krzywizna) wykracza poza Evesa -- Eves \cite[s. 315]{eves1_1972} sam jej nie rozwija; potrzebne inne źródło, jeśli chcemy ją opisać.} % lang-pl
\todofoot{La geometria riemanniana (metrica, curvatura) va oltre Eves -- Eves \cite[p. 315]{eves1_1972} non la sviluppa; serve un'altra fonte, se vogliamo descriverla.} % lang-it
